\chapter{Introducción}

En la Introducción se debe resumir de forma esquemática pero suficientemente clara lo esencial de cada una de las partes del trabajo y, sobre todo, aprovechar para explicar conceptos que puedan ser de utilidad después. 
\newline

La lectura de esta parte debe contextualizar perfectamente todo el Trabajo y debe estar PLAGADA DE REFERENCIAS (más adelante se describe cómo hacerlo).
\newline

Es una parte muy importante de la memoria. Las ideas principales a transmitir son la identificación del problema a tratar, la justificación de su importancia, esbozo de los objetivos generales a grandes rasgos y un adelanto de la contribución que esperas hacer.

A modo de guía, la Introducción debe contener estos tres apartados:
\begin{itemize}
\item Motivación / justificación del tema a tratar
\item Planteamiento del Trabajo (atención a la mayúscula de "Trabajo": cuando nos referimos a nuestro documento, ese Trabajo va en mayúscula)
\item Estructura del Trabajo
\end{itemize}

\vspace{1cm}
{\bf{ATENCIÓN, CÓMO REFERENCIAR.}}

A continuación vamos a ver qué son, qué no son y para qué  sirven las referencias bibliográficas:
\begin{itemize}
	\item las referencias NO están para rellenar
	\item las referencias SÍ son un TRIBUTO a las personas que hicieron en primer lugar una investigación o aportaron una idea
	\item la referencias citan trabajo que acompaña una idea o un concepto, NO se usa UNA referencia para todo un párrafo que aglutina diferentes conceptos
	\item deben citarse LOS TRABAJOS ORIGINALES DE LOS AUTORES, y no un libro de texto donde he visto que hablan de algo
	\item os libros de texto se pueden referencial excepcionalmente, porque el libro es sobresaliente y conocido en el campo y/o porque quiero apoyarme en una explicación conocida para un determinado concepto o cuestión, pero no como origen del concepto del que se habla
\end{itemize}

\vspace{1cm}
{\bf{VEAMOS EJEMPLOS DE CITAS.}}

Vamos a ver los dos tipos de citas que vamos a utilizar normalmente en los trabajos y actividades:
\begin{itemize}
	\item CITA INDIRECTA: estamos hablando de algo e introcimos la referencia sobre la que apoyamos esa idea. Por ejemplo, hablado de los procesadores de texto, LATEX  \citep{lamport1994} es tipo de editor con el que se consiguen resultados de gran calidad
	\item CITA DIRECTA: nos referimos directamente a un trabajo y lo citamos. Por ejemplo, según \cite{ackerman2017}, la liga de fútbol inglesa debe tener torneos de desempate
\end{itemize}

Según el párrafo anterior con la cita indirecta aparecerá entre paréntesis el apellido y el año, separados por una coma; en el caso de varios autores, podrán aparecer hasta tres apellidos y, cuando son más, aparecerá el primero seguido de "et al." (y después el año separado con una coma). Con la cita directa solo va entre paréntesis el año, pero no los autores. Podéis mirar el apartado de Bibliografía para ver cómo han quedado estas citas.
\newline

En esta Actividad, cuando vamos a resolver problemas, PUEDE utilizarse libros de texto como referencias de procedimientos matemáticos que vamos a utilizar, pero sin olvidar el párrafo anterior. Es decir, sería conveniente citar ARTÍCULOS ORIGINALES de donde surgieron esos desarrollos.
\newline

Para consultar detenidamente la formas de referenciar, debéis seguir el formato APA del siguiente documento: \url{https://bibliografiaycitas.unir.net/documentos/APA_7ed_UNIR.pdf}